\section{Endogenous common decisions}\label{sec:v13-controlled}
\subsection{A shared encoder on a controlled tree}
Let $\mathcal H_t$ be the finite supported observation--control histories
of a known controlled experiment. The initial history has distribution
$\mu$, positive on $\mathcal H_0$. For $h\in\mathcal H_t$ and $a\in A_t$,
write $p_t(y\mid h,a)$ for the next-report law and $hay$ for the successor
when this probability is positive. Zero-probability successors impose no
condition. The finite task set is $D$. Task $d$ has conditional expected
stage costs $c_{t,d}(h,a)$ and terminal cost $c_{T,d}(h)$. These may be
obtained by conditioning bounded losses of the actual hidden mark;
the task does not alter the physical transition law. Controls do alter
that law through $a$.

The common source-consuming encoder is
\[
 q_0(m\mid h),\qquad
 q_{t+1}(m'\mid m,a,y),\quad m\in[S_t],\quad m'\in[S_{t+1}].
\]
Controller $d$ chooses its action by $\alpha_{t,d}(a\mid m)$. It cannot
query the discarded history. The clock is available. The encoder does
not know $d$; the update sees the actual selected action, as allowed by
the interface. All kernels can be stochastic. The same encoder must
serve every task, even though different tasks induce different laws
on histories.

Define the full-history Bellman functions and regrets by
\begin{align}
 V_{T,d}(h)&=c_{T,d}(h),\nonumber\\
 V_{t,d}(h)&=\min_{a\in A_t}
   \left(c_{t,d}(h,a)+\sum_y p_t(y\mid h,a)V_{t+1,d}(hay)\right),
                                      \label{eq:v13-bellman}\\
 G_{t,d}(h,a)&=c_{t,d}(h,a)+\sum_y p_t(y\mid h,a)V_{t+1,d}(hay)
                                  -V_{t,d}(h)\ge0.\nonumber
\end{align}
The benchmark for task $d$ is $v_d=\sum_h\mu(h)V_{0,d}(h)$.
It is an ordinary full-history control value, not a posterior quantity
conditioned on a fixed open-loop preparation.

\begin{lemma}[Occupation regret]\label{lem:v13-regret}
For every admissible common encoder and task controller, the task excess
satisfies the exact identity
\begin{equation}\label{eq:v13-regret}
 R_d-v_d=\E_d\sum_{t=0}^{T-1}G_{t,d}(H_t,A_t).
\end{equation}
Moreover the entire vector of risks of a randomized finite-horizon
design is a convex combination of risk vectors of deterministic designs
with the same widths and one common encoder.
\end{lemma}
\begin{proof}
Take conditional expectations in \eqref{eq:v13-bellman} along the actual
controlled path. Summing the resulting expressions for $G_{t,d}$
telescopes the terms $V_{t,d}(H_t)$, leaving the sum of stage and terminal
costs minus $\E V_{0,d}(H_0)$. This proves \eqref{eq:v13-regret} without
assuming the control is open loop.

Sample in advance a deterministic outcome for every time-indexed row
of the common encoder and every time-indexed row of every task
controller. Row choices are independent with their prescribed
probabilities; the sampled encoder table is shared across tasks.
Along each trajectory each time-indexed row is encountered at most
once. The product of its row probabilities is therefore the probability
of that path under the sampled-table mixture. This identity holds for
all tasks simultaneously after summing paths and costs. An additional
independent visible or hidden seed is first conditioned on and then
integrated. The conclusion is a mixture identity; the mixing variable
is not claimed to be executable privately without retaining its index.
\end{proof}

\begin{definition}[Endogenous compatible cover]\label{def:v13-cover}
A cover of profile $S$ consists of an initial map
$i:\mathcal H_0\to[S_0]$, common maps
$\delta_t:[S_t]\times A_t\times Y_{t+1}\to[S_{t+1}]$,
task action maps $a_{t,d}:[S_t]\to A_t$, and cells
$C_{t,m}^d\subseteq\mathcal H_t$. Empty cells and overlapping cells are
allowed. The requirements are
\begin{align}
 h&\in C_{0,i(h)}^d&& (h\in\mathcal H_0,\ d\in D),\label{eq:v13-cover}\\
 G_{t,d}(h,a_{t,d}(m))&=0&& (h\in C_{t,m}^d),\nonumber\\
 h\in C_{t,m}^d,\ p_t(y\mid h,a_{t,d}(m))>0
 &\ \Longrightarrow\
 ha_{t,d}(m)y\in C_{t+1,\delta_t(m,a_{t,d}(m),y)}^d.\nonumber
\end{align}
Thus the cells cover the subtree generated by the selected actions.
They need not cover histories generated by actions the controller never
selects. Crucially, $i$ and $\delta_t$ cannot depend on $d$.
\end{definition}

\begin{theorem}[Exact endogenous compression]\label{thm:v13-endogenous}
For a prescribed finite width profile $S$, the following are equivalent:
a deterministic common design attains $R_d=v_d$ for every task;
a privately randomized common design does so;
a common design with an independent visible or hidden seed does so;
and an endogenous compatible cover of profile $S$ exists.

Let $\beta_d>0$ sum to one, and put
$\mathcal G=\sum_d\beta_d(R_d-v_d)$. If the profile is infeasible,
then every such design satisfies
\begin{equation}\label{eq:v13-gap}
 \mathcal G\ge g_S:=\beta_*\mu_*p_*^T\gamma_*>0,
\end{equation}
where $\beta_*=\min_d\beta_d$, $\mu_*$ is the least positive initial
mass, $p_*$ is the least positive transition probability, and
$\gamma_*$ is the least positive Bellman regret among supported entries.
If there are no positive regrets, every profile with positive widths is
feasible. The convention $p_*=1$ applies when there are no transitions.

More generally, replacing the zero in \eqref{eq:v13-cover} by the
requirement $G_{t,d}(h,a_{t,d}(m))\le\eta_{t,d}$ gives a common design
with $R_d-v_d\le\sum_t\eta_{t,d}$. The sharper value is exactly the
occupation-weighted sum in \eqref{eq:v13-regret}.
\end{theorem}
\begin{proof}
A cover defines the deterministic encoder $M_0=i(H_0)$ and
$M_{t+1}=\delta_t(M_t,A_t,Y_{t+1})$, and the task controller uses
$a_{t,d}(M_t)$. Induction in \eqref{eq:v13-cover} shows that every reached
history lies in its indicated cell. Each incurred Bellman regret is
zero, so Lemma~\ref{lem:v13-regret} proves exactness. The same induction
and identity prove the approximate-cover statement.

Conversely, an exact deterministic design defines $C_{t,m}^d$ to be the
set of histories reached with positive probability under task $d$ and
state $m$. Initial coverage and common successor closure follow from
its encoder. Since the nonnegative sum in \eqref{eq:v13-regret} has
zero expectation, every positively reached history has zero regret
for its selected action. This proves all cover requirements. It also
explains why adding closure for every unused action would impose a
strictly different problem.

For a randomized design, Lemma~\ref{lem:v13-regret} expresses the vector
of excesses as an average of deterministic excess vectors. Each
coordinate is nonnegative. If all averaged coordinates are zero, then,
for almost every sampled design, all coordinates are zero, since the
task set is finite. Choose one such deterministic table. This applies
after conditioning on an arbitrary independent seed, with recovery
from the seed and retained state when the seed is visible. It does not
use a seed correlated with the initial history. The deterministic
cover already constructed then proves necessity.

If no exact deterministic design exists, each deterministic design has
some task, stage, and positively reached history with positive selected
regret. That history has probability at least $\mu_*p_*^t$, and its
regret is at least $\gamma_*$. Since $0<p_*\le1$ and $t\le T$, its
contribution to $\mathcal G$ is at least the right side of
\eqref{eq:v13-gap}. Every randomized design is an average of these
vectors, so the same bound follows. When there is no positive regret,
every action is Bellman optimal wherever supported, and a one-state
encoder with arbitrary task actions is exact.
\end{proof}

\begin{example}[Why unused controls cannot define the cover]
At the first time there is no report. Action $\mathtt{rest}$ leads to
a deterministic hidden bit zero and a constant report, at zero cost.
Action $\mathtt{probe}$ costs one, draws a fair bit, and reports it.
After this report is consumed there is one blank checkpoint, followed
by a binary guess with zero--one loss. The full-history optimum chooses
$\mathtt{rest}$ and guesses zero; one retained state is exact. A restart
family that insists on optimal guessing after the unused
$\mathtt{probe}$ branch needs two retained states. The endogenous
criterion and the all-controls restart criterion are both meaningful,
but they have different feasible sets. The cover in
Definition~\ref{def:v13-cover} implements the former, not the latter.
\end{example}

\subsection{Lossy common design and intrinsic comparison}
Let $\mathcal D_S^{\rm ctl}$ be the finite set of deterministic common
encoders and deterministic task controllers in this section. Put
$r_d(\pi)=R_d(\pi)-v_d$. Then the hidden-selector lossy minimax value is
\begin{equation}\label{eq:v13-lossy}
 \min_{\lambda\in\Delta(\mathcal D_S^{\rm ctl})}
          \max_d\sum_\pi\lambda_\pi r_d(\pi)
 =\max_{\beta\in\Delta(D)}\min_{\pi\in\mathcal D_S^{\rm ctl}}
          \sum_d\beta_dr_d(\pi).
\end{equation}
Here a selector chooses the entire design, including its task-controller
tables; its mode is internal to that design. This follows from the
simultaneous mixture identity in Lemma~\ref{lem:v13-regret} and finite
minimax. A mixture using at most
$|D|+1$ designs suffices and can be stored with at most $(|D|+1)S_t$
private labels. Private finite-width optimization is instead the
polynomial optimization in the actual encoder and controller rows.
The exact zero sets agree by Theorem~\ref{thm:v13-endogenous}, but
positive minimax values need not agree. Thus the cover theorem supplies
a distinguished zero face inside a quantitative lossy theory, rather
than purporting to solve every lossy design problem by equivalence
classes.

\begin{corollary}[Intrinsic certification of endogenous widths]
\label{cor:v13-certificate}
Suppose the losses of this finite problem are normalized so that every
total task loss lies in $[0,1]$. Let an approximation $F$ have the same
control and task interfaces, with each task represented by the same
bounded marked-transcript loss on both sides and the full-history
benchmark included in the paired decision class. If
\[
 \max\{\dpr_{\mathbf1}(E,F),\dpr_{\mathbf1}(F,E)\}
                    \le\varepsilon<g_S/4,
\]
then the optimized $\beta$-weighted excess at profile $S$ in $F$ is
less than $g_S/2$ exactly when the endogenous cover exists in $E$.
If only a cost-$K$ simulation from $E$ to $F$ is known and an exact
$S$-state $F$ design transports with zero error and the same benchmark,
then an endogenous cover of profile $KS$ exists in $E$. A hidden
selector multiplies that profile by its charged selector size.
\end{corollary}
\begin{proof}
Full-history controllers also transport through a stateless simulator,
so their optimized benchmarks differ by at most $\varepsilon$.
Equation~\eqref{eq:v13-risktransport} bounds the optimized constrained
values by $\varepsilon$ in each direction. The optimized excesses
therefore differ by at most $2\varepsilon$. In $E$ the value is zero
on the cover face and at least $g_S$ otherwise, by
Theorem~\ref{thm:v13-endogenous}. The strict threshold follows. For the
one-sided statement, the product construction transports the exact
design at the stated width and benchmark. The same theorem then
converts exactness into a cover.
\end{proof}

\begin{remark}[Dependence of the finite witness]
The number $g_S$ depends on the least positive probability and regret
of the specified finite model. It can tend to zero as a transition or
a cost is perturbed. Neither \eqref{eq:v13-gap} nor the certification
corollary asserts a model-uniform robustness radius. Likewise an
inverse-limit assertion that every finite prefix has a feasible cover
would give existence of an infinite consistent design, not a computable
bound on the first infeasible prefix. No such bound is used here.
\end{remark}
